\documentclass[11pt]{article}

\usepackage[utf8]{inputenc}
\usepackage[T1]{fontenc}
\usepackage{geometry}
\usepackage{amsmath}
\usepackage{booktabs}
\usepackage{array}

\geometry{letterpaper, margin=2.5cm}

\title{Informe del módulo de búsqueda web}
\author{CSDI RAG Engine}
\date{}

\begin{document}

\maketitle

\section{Módulo de búsqueda web}

El módulo de búsqueda web es el mecanismo que permite al sistema responder consultas sobre
información que no está presente en el corpus indexado. Su responsabilidad no es sustituir la
recuperación local, sino complementarla cuando el corpus no ofrece suficiente evidencia para
responder la pregunta del usuario. El módulo se activa automáticamente, de forma transparente
para el usuario, como parte del pipeline RAG.

Para que esta activación sea selectiva y no sistemática, el módulo incorpora un componente
dedicado a medir la calidad de los resultados locales antes de recurrir a la web. Solo cuando
esa calidad cae por debajo de un umbral configurable el sistema ejecuta la búsqueda externa.
Esto preserva la rapidez del pipeline en el caso común, donde el corpus contiene la información
necesaria, y reserva el costo de la búsqueda web para los casos en que realmente añade valor.

\paragraph{Arquitectura del módulo.}
El módulo de búsqueda web se organiza en cuatro componentes con responsabilidades distintas:
el \emph{detector de insuficiencia}, que decide si la evidencia local es suficiente; el
\emph{orquestador de búsqueda web}, que coordina la búsqueda externa cuando el detector lo
requiere; el \emph{proveedor de búsqueda}, que ejecuta la consulta contra un motor externo; y
el \emph{descargador de documentos}, que obtiene y limpia el contenido de las páginas
encontradas. Un quinto componente transversal es el \emph{índice de caché web}, que persiste
los documentos descargados para que puedan ser recuperados en consultas futuras sin repetir la
búsqueda externa.

Estos componentes no operan en paralelo: forman una cadena activada condicionalmente dentro
del pipeline RAG. La cadena completa solo se ejecuta cuando el corpus local no satisface la
consulta, y dentro de ella se intenta primero el caché antes de invocar el motor externo.

\paragraph{Diseño general del flujo de búsqueda web.}
Cuando el pipeline RAG recupera fragmentos del corpus y los entrega al detector de
insuficiencia, el detector calcula una puntuación de confianza local. Si esa puntuación supera
el umbral configurado, el pipeline procede directamente a la generación de respuesta. Si no lo
supera, el pipeline consulta primero el caché web: fragmentos de búsquedas anteriores ya
indexados. Si el caché tampoco es suficiente, se ejecuta la búsqueda externa, se descargan y
chunkean los documentos encontrados, se indexan en el caché web y se recuperan inmediatamente
para la consulta actual. En cualquier punto donde la suficiencia se recupere, la cadena se
detiene.

\section{Detector de insuficiencia}

El detector de insuficiencia (\texttt{InsufficiencyDetector}) es el componente que determina
si los fragmentos recuperados del corpus local son suficientes para responder la consulta. Esta
decisión no se toma con un umbral simple sobre el número de resultados ni sobre la puntuación
máxima: se calcula una puntuación de confianza compuesta a partir de cinco señales
independientes que capturan aspectos distintos de la calidad de la recuperación.

\paragraph{Métricas calculadas.}
El detector calcula nueve valores sobre el conjunto de fragmentos recuperados, de los cuales
cinco contribuyen directamente a la puntuación de confianza final:

\begin{itemize}
\item \textbf{top\_score\_norm}: puntuación máxima de los fragmentos, normalizada respecto al
máximo teórico del método de fusión. Cuando la fusión es RRF con constante $k$ y suma de pesos
$W$, el máximo teórico es $W/(k+1)$; con $k=60$ y $W=1{,}0$ ese máximo es
$1/61 \approx 0{,}0164$. Normalizar permite comparar el score de forma independiente a la
escala de los recuperadores.

\item \textbf{quantity\_score}: proporción de fragmentos recuperados respecto a la cantidad
esperada, calculada como $\min(1{,}0,\, n_{\text{results}} / n_{\text{expected}})$ con
$n_{\text{expected}}=10$. Satura en 1 cuando hay al menos diez fragmentos.

\item \textbf{coverage\_score}: fracción promedio de términos de la consulta presentes en los
$N$ fragmentos de mayor puntuación (por defecto $N=5$). Para cada fragmento se calcula el
cociente entre los términos de la consulta que aparecen en su texto y el total de términos de
la consulta; el coverage\_score es la media de esos cocientes. Mide si los fragmentos
recuperados contienen el vocabulario de la pregunta.

\item \textbf{diversity\_score}: cociente entre el número de URLs únicas y el número total
de fragmentos. Un valor bajo indica que los fragmentos provienen de una sola fuente, lo que
reduce la robustez de la respuesta frente a errores o sesgos de esa fuente.

\item \textbf{answerability\_score}: estimación de la capacidad de respuesta combinando dos
señales: la proporción de fragmentos con cobertura suficiente de la consulta
(\emph{relevant\_ratio}) y el coverage\_score general. La fórmula es:
\[
\text{answerability} = \mathrm{clamp}_{[0,1]}\!\left(0{,}6 \cdot
\frac{\min(r,\, r_{\min})}{r_{\min}} + 0{,}4 \cdot \text{coverage\_score}\right)
\]
donde $r$ es el número de fragmentos cuyo overlap con la consulta supera el umbral
$\theta_{\text{overlap}}=0{,}20$ y $r_{\min}=2$ es el mínimo esperado. El coeficiente 0,6
pondera la profundidad (¿hay al menos dos fragmentos que hablan del tema?) y el 0,4 la amplitud
(¿cubren el vocabulario de la consulta?).
\end{itemize}

\paragraph{Tokenización para cobertura y answerability.}
El cálculo de cobertura requiere comparar los términos de la consulta con los del texto de cada
fragmento. El tokenizador usado, \texttt{simple\_tokenize}, aplica expresiones regulares
Unicode (\texttt{[\textasciicircum\textbackslash W\_]+}) para extraer palabras en minúscula,
filtra stopwords en inglés y en español, y descarta tokens de un solo carácter. El soporte
explícito de stopwords en español garantiza que la cobertura no se infle por palabras vacías en
consultas formuladas en ese idioma.

\paragraph{Fórmula de confianza local.}
Las cinco métricas se combinan en una puntuación de confianza ponderada:
\[
c_{\text{local}} = \mathrm{clamp}_{[0,1]}\!\left(
  w_{\text{top}} \cdot \text{top\_score\_norm}
  + w_q \cdot \text{quantity\_score}
  + w_c \cdot \text{coverage\_score}
  + w_d \cdot \text{diversity\_score}
  + w_a \cdot \text{answerability\_score}
\right)
\]
con valores por defecto $w_{\text{top}}=0{,}10$, $w_q=0{,}15$, $w_c=0{,}35$,
$w_d=0{,}15$, $w_a=0{,}25$. Los pesos suman exactamente 1,0; el sistema valida esta
invariante al arrancar y rechaza configuraciones inválidas.

\paragraph{Justificación de los pesos.}
La distribución de pesos no es uniforme porque las métricas no tienen el mismo poder predictivo
sobre la calidad de la recuperación. El coverage\_score recibe el mayor peso (0,35) porque es
el mejor proxy disponible de relevancia léxica sin necesitar un modelo adicional: si los
fragmentos recuperados no contienen los términos de la consulta, es improbable que puedan
responderla. El answerability\_score recibe el segundo peso más alto (0,25) porque combina
dos señales y captura tanto la profundidad como la anchura de la cobertura. El top\_score\_norm
recibe el peso más bajo (0,10) deliberadamente: una puntuación alta de fusión indica que el
recuperador híbrido encontró buenos candidatos según sus propias métricas de similitud, pero no
garantiza que esos fragmentos contengan la respuesta a la pregunta específica del usuario.

\paragraph{Decisión de búsqueda web y razones diagnósticas.}
La decisión de activar la búsqueda web se toma comparando $c_{\text{local}}$ con el umbral
$\theta_c=0{,}65$: si $c_{\text{local}} < \theta_c$, el sistema requiere búsqueda web. De
forma independiente a ese umbral, el detector evalúa condiciones individuales sobre cada
métrica y construye una lista de razones diagnósticas. Las razones posibles son:
\texttt{NO\_RESULTS} (sin fragmentos), \texttt{LOW\_NUM\_RESULTS} (menos de 5 fragmentos),
\texttt{LOW\_TOP\_SCORE} (top\_score\_norm $< 0{,}35$), \texttt{LOW\_COVERAGE} (coverage $<
0{,}20$), \texttt{LOW\_SOURCE\_DIVERSITY} (diversity $< 0{,}30$), \texttt{LOW\_ANSWERABILITY}
(answerability $< 0{,}40$) y \texttt{LOW\_CONFIDENCE} (confianza global bajo el umbral). Esta
lista de razones se devuelve junto con la decisión y permite auditar por qué se activó la
búsqueda en cada consulta, lo que facilita el diagnóstico y el ajuste de parámetros.

\section{Caché web}

Antes de ejecutar una búsqueda externa, el pipeline consulta el caché web: un índice híbrido
independiente del corpus principal que almacena los fragmentos obtenidos en búsquedas web
anteriores. La separación física entre el corpus principal y el caché web es una decisión de
diseño deliberada: garantiza que los documentos descargados de la web no contaminen el índice
del corpus curado y que puedan gestionarse con políticas distintas de retención y actualización.

\paragraph{Estructura del caché web.}
El caché web replica la arquitectura de índice del corpus principal pero en tablas separadas:
\texttt{web\_cache\_chunks} para los fragmentos de texto, \texttt{web\_cache\_bm25\_*} para el
índice invertido léxico y \texttt{web\_cache\_vector\_*} para el índice FAISS de embeddings.
La misma clase \texttt{HybridRetriever} que opera sobre el corpus puede operar sobre el caché
web con distintas instancias de repositorios, reutilizando toda la lógica de fusión RRF y
reordenamiento sin duplicar código.

Los documentos completos descargados se persisten en \texttt{web\_cache\_documents} con
deduplicación por URL mediante un \texttt{INSERT ... ON CONFLICT (url) DO UPDATE}. Esto evita
indexar el mismo documento dos veces si la misma URL aparece en búsquedas distintas.

\paragraph{Evaluación del caché antes de la búsqueda externa.}
Cuando el detector de insuficiencia indica que el corpus no es suficiente, el pipeline ejecuta
una recuperación sobre el caché web y somete los resultados al mismo detector. Si los fragmentos
del caché elevan la confianza por encima del umbral, la búsqueda externa no se ejecuta. Este
mecanismo de doble evaluación convierte el caché en un amortiguador eficiente: búsquedas
repetidas sobre temas similares reutilizan los documentos ya descargados sin incurrir en el
costo de una nueva llamada al motor externo.

\section{Orquestador de búsqueda web}

El \texttt{WebSearchOrchestrator} es el componente que coordina la búsqueda externa cuando
ni el corpus ni el caché son suficientes. Recibe la consulta, la pasa al proveedor de búsqueda,
descarga los documentos encontrados, los fragmenta, los indexa en el caché web y persiste un
registro de la búsqueda. Todo esto ocurre dentro de una sola llamada a \texttt{run(query)},
que devuelve un \texttt{WebSearchRunResult} con los hits obtenidos, los documentos descargados
y el número de fragmentos indexados.

Si la búsqueda está deshabilitada (\texttt{WEB\_SEARCH\_ENABLED=false}) o la consulta está
vacía, el orquestador devuelve un resultado vacío sin ejecutar ninguna operación. Este punto de
desactivación explícita permite operar el sistema en modo \emph{corpus-only} durante pruebas o
en entornos sin acceso a internet.

\paragraph{Persistencia de búsquedas.}
Cada ejecución del orquestador se registra en la tabla \texttt{web\_search\_runs} con la
consulta normalizada, el proveedor utilizado, el número de hits obtenidos y el número de
fragmentos indexados. Este registro permite auditar la actividad de búsqueda externa, detectar
patrones de insuficiencia recurrentes en el corpus y evaluar la efectividad del caché a lo
largo del tiempo.

\section{Proveedor de búsqueda: DuckDuckGo}

El proveedor de búsqueda implementado es \texttt{DuckDuckGoSearchProvider}, que consulta la
interfaz HTML de DuckDuckGo sin requerir clave de API. La búsqueda se realiza mediante una
petición HTTP GET al endpoint \texttt{https://duckduckgo.com/html/} con el parámetro
\texttt{q=<consulta>}, y los resultados se extraen del HTML devuelto mediante selectores CSS
con BeautifulSoup.

\paragraph{Normalización de enlaces.}
DuckDuckGo encapsula las URLs de los resultados en redirecciones internas con el formato
\texttt{/l/?uddg=<url\_codificada>}. El proveedor detecta este patrón y decodifica la URL
real antes de devolver el hit, evitando que el descargador intente acceder a la URL de
redirección en lugar del documento original.

\paragraph{Justificación de la elección del proveedor.}
DuckDuckGo fue seleccionado porque no requiere registro ni clave de API para uso programático
de su interfaz HTML, lo que elimina una dependencia de configuración externa para el
funcionamiento básico del sistema. La contrapartida es que la interfaz HTML no es una API
oficial y puede cambiar sin aviso, lo cual se identificó explícitamente como una limitación
del diseño (véase sección de limitaciones). Para entornos de producción con mayor estabilidad,
el diseño de la interfaz \texttt{SearchProvider} como \texttt{Protocol} permite sustituir el
proveedor por cualquier implementación alternativa sin modificar el resto del módulo.

\section{Descargador de documentos}

El \texttt{HttpDocumentFetcher} descarga el contenido de cada URL retornada por el proveedor.
Para cada hit, realiza una petición HTTP GET con un \emph{User-Agent} identificable, verifica
el tipo de contenido de la respuesta y selecciona la estrategia de extracción de texto
apropiada.

\paragraph{Extracción de contenido HTML.}
Para respuestas con \texttt{Content-Type: text/html}, el descargador delega en el
\texttt{Scraper} del sistema, que aplica selectores configurados para extraer el contenido
principal de la página: se seleccionan elementos \texttt{main}, \texttt{article} o
\texttt{[role=main]}, excluyendo \texttt{script}, \texttt{style} y \texttt{noscript}. El
título se extrae del primer \texttt{h1} o del elemento \texttt{title}. Para respuestas no
HTML, se usa el texto plano de la respuesta directamente.

Tras la extracción, el texto se normaliza colapsando secuencias de espacios y saltos de línea
en un único espacio, y se trunca a un máximo de 10\,000 caracteres. Este límite evita que
páginas muy largas generen un número excesivo de fragmentos y degraden el rendimiento del
índice. Si el texto resultante está vacío, el descargador devuelve \texttt{None} y el
orquestador omite ese documento. Los fallos de red o HTTP se registran en el log y no
interrumpen la descarga de los demás documentos del mismo lote.

\paragraph{Fragmentación e indexación.}
Los documentos descargados se fragmentan con el mismo \texttt{Chunker} usado para el corpus
principal (256 palabras, solapamiento de 32 palabras). Cada fragmento recibe un
\texttt{source\_id} con el formato \texttt{web:<proveedor>:<dominio>}, por ejemplo
\texttt{web:duckduckgo:docs.python.org}, y un campo \texttt{breadcrumb="web-search"}.
Este campo permite que el pipeline identifique de forma unívoca qué fragmentos de la respuesta
final provienen de búsquedas web frente a los que provienen del corpus curado.

Los fragmentos se ingresan en el caché web mediante el mismo servicio de ingesta del corpus,
que actualiza tanto el índice BM25 como el vectorial. Tras la ingesta, el índice BM25 del
caché se reconstruye para reflejar los nuevos documentos antes de que el pipeline realice la
recuperación post-búsqueda.

\section{Integración en el pipeline RAG}

El módulo de búsqueda web no opera de forma independiente: es una rama condicional del
\texttt{RAGPipeline}. La decisión de activarla, el uso del caché y la recuperación
post-búsqueda están coordinados por el mismo pipeline que gestiona la recuperación del corpus
y la generación de respuesta.

\paragraph{Flujo condicional completo.}
El pipeline evalúa la suficiencia en dos momentos: después de la recuperación del corpus y
después de la recuperación del caché web. El flujo completo es:

\begin{enumerate}
\item Recuperación híbrida sobre el corpus principal.
\item Evaluación de suficiencia con \texttt{InsufficiencyDetector}.
\item Si insuficiente y caché habilitado: recuperación sobre el caché web y segunda evaluación.
\item Si aún insuficiente y búsqueda externa habilitada: ejecución del orquestador, indexación
      de documentos nuevos y recuperación sobre el caché actualizado.
\item Reordenamiento con cross-encoder sobre el conjunto final de fragmentos.
\item Generación de respuesta con el LLM.
\end{enumerate}

En cualquier punto donde la suficiencia se recupere, los pasos siguientes no se ejecutan. Si
el \texttt{WebSearchOrchestrator} no está configurado (instancia \texttt{None}), el pipeline
responde con los fragmentos disponibles del corpus sin intentar la búsqueda externa, lo que
permite operar el sistema en modo degradado sin errores.

\paragraph{Trazabilidad de fuentes en la respuesta.}
El resultado devuelto al cliente incluye una lista de fuentes con el campo
\texttt{source\_type}, que distingue entre \texttt{"corpus"} y \texttt{"web\_cache"}. Esta
distinción permite al usuario saber si la información proviene del corpus curado del sistema o
de documentos obtenidos en tiempo real de la web, lo que es relevante para evaluar la
fiabilidad de la respuesta. Adicionalmente, los campos \texttt{external\_search\_executed} y
\texttt{external\_indexed\_count} del resultado informan si se ejecutó una búsqueda externa y
cuántos fragmentos se indexaron.

\section{Configuración y parámetros clave}

El módulo de búsqueda web expone 23 parámetros configurables mediante variables de entorno,
organizados en grupos funcionales. Los más relevantes para el comportamiento del sistema son:

\begin{center}
\begin{tabular}{lll}
\toprule
\textbf{Variable} & \textbf{Default} & \textbf{Efecto} \\
\midrule
\texttt{WEB\_SEARCH\_ENABLED} & true & Activa o desactiva el módulo completo \\
\texttt{WEB\_SEARCH\_TOP\_K} & 5 & Resultados solicitados al proveedor \\
\texttt{WEB\_CACHE\_ENABLED} & true & Consulta el caché antes de búsqueda externa \\
\texttt{WEB\_CACHE\_TOP\_K} & 10 & Fragmentos recuperados del caché \\
\texttt{INSUFF\_CONFIDENCE\_THRESHOLD} & 0.65 & Umbral global de confianza \\
\texttt{INSUFF\_W\_COVERAGE} & 0.35 & Peso del coverage en la confianza \\
\texttt{INSUFF\_W\_ANSWERABILITY} & 0.25 & Peso del answerability en la confianza \\
\texttt{INSUFF\_W\_TOP} & 0.10 & Peso del top\_score en la confianza \\
\texttt{INSUFF\_MIN\_COVERAGE\_SCORE} & 0.20 & Mínimo de coverage para no marcar razón \\
\texttt{INSUFF\_MIN\_ANSWERABILITY\_SCORE} & 0.40 & Mínimo de answerability \\
\texttt{INSUFF\_COVERAGE\_TOP\_N} & 5 & Fragmentos evaluados para coverage \\
\bottomrule
\end{tabular}
\end{center}

La suma de los cinco pesos del detector es validada al cargar la configuración: si no suman
exactamente 1,0 (con tolerancia $10^{-6}$), el sistema rechaza la configuración con un error
explícito. Esto evita que configuraciones incorrectas produzcan puntuaciones de confianza fuera
del rango $[0,1]$ de forma silenciosa.

\section{Justificación técnica de la solución}

El módulo de búsqueda web existe porque ningún corpus estático puede anticipar todas las
consultas posibles de los usuarios. Documentación desactualizada, temas de larga cola no
cubiertos por el corpus, o preguntas sobre eventos recientes son casos que la recuperación
local no puede resolver por diseño. Sin un mecanismo de extensión, el sistema generaría
respuestas incorrectas o reconocería su ignorancia en situaciones donde la información existe
y es accesible.

La decisión de no activar la búsqueda web en cada consulta responde a un balance entre calidad
y costo. Una búsqueda externa en cada consulta añadiría latencia constante, dependencia de un
servicio externo y ruido potencial de documentos no relevantes, sin beneficio cuando el corpus
ya contiene la respuesta. El detector de insuficiencia resuelve este balance estimando la
calidad local antes de incurrir en ese costo.

\paragraph{Por qué una puntuación compuesta y no un umbral simple.}
Un umbral simple sobre el número de resultados o sobre la puntuación máxima del recuperador
sería insuficiente porque cada señal captura solo una dimensión del problema. Una consulta
puede tener muchos resultados con puntuaciones altas pero todos provenientes de la misma fuente
y sin cubrir los términos de la pregunta: el número y el score serían altos pero la respuesta
sería de baja calidad. La puntuación compuesta detecta este caso porque el coverage\_score y
el diversity\_score serían bajos, arrastrando la confianza por debajo del umbral.

\paragraph{Por qué el caché web es un componente separado y no parte del corpus.}
Los documentos del corpus son curados y estructurados; los documentos de la web son
heterogéneos, potencialmente ruidosos y temporalmente válidos. Mezclarlos en el mismo índice
degradaría la calidad del corpus y complicaría las operaciones de mantenimiento como la
reconstrucción del índice o la actualización de documentos. La separación permite aplicar
políticas distintas a cada índice y garantiza que el corpus principal no se vea afectado por
el volumen o calidad variable de los documentos web.

\paragraph{Por qué se reutiliza la arquitectura de recuperación del corpus.}
El caché web usa las mismas clases \texttt{HybridRetriever}, \texttt{BM25Retriever} y
\texttt{VectorRetriever} que el corpus, instanciadas con repositorios distintos. Esta decisión
evita duplicar la lógica de recuperación y garantiza que el caché se beneficie de las mismas
mejoras que se apliquen al corpus: si se ajustan los pesos de fusión o se cambia el modelo de
embeddings, el caché hereda esos cambios sin intervención adicional.

\section{Limitaciones del módulo}

El módulo de búsqueda web tiene limitaciones conocidas que acotan su alcance y confiabilidad.

El proveedor DuckDuckGo opera sobre la interfaz HTML pública, que no es una API oficial. Esta
interfaz puede cambiar su estructura sin aviso previo, lo que puede romper el parser de
resultados. En ese caso, el módulo devolvería cero resultados sin error explícito hasta que
el extractor sea actualizado.

La descarga de documentos falla silenciosamente para URLs que requieren autenticación, que
devuelven contenido dinámico generado por JavaScript, o que bloquean el \emph{User-Agent}
del sistema. En todos esos casos el documento es omitido, lo que puede reducir el número
efectivo de fragmentos disponibles tras la búsqueda.

El truncado de documentos a 10\,000 caracteres puede cortar información relevante en páginas
largas. Para documentación técnica extensa, como páginas de referencia de API o tutoriales
completos, los fragmentos finales del documento pueden quedar fuera del límite y no ser
indexados.

El caché web no tiene política de expiración implementada. Documentos indexados en búsquedas
anteriores permanecen en el caché indefinidamente, lo que puede provocar que consultas
posteriores recuperen información desactualizada sin que el sistema lo detecte. Este problema
es especialmente relevante para consultas sobre información que cambia con el tiempo.

La evaluación de suficiencia se basa en señales léxicas (coverage) y de cantidad, no en una
estimación semántica profunda de si los fragmentos contienen la respuesta correcta. Es posible
que el detector no active la búsqueda web en casos donde los fragmentos recuperados son
temáticamente relevantes pero factualmente incorrectos o incompletos.

\end{document}
